\begin{table}[H]
\centering
\caption*{\textbf{Quadro A.2 - Contratos de execução, evidência e revisão}}
\scriptsize
\renewcommand{\arraystretch}{1.12}
\begin{tabularx}{\linewidth}{|p{0.18\linewidth}|p{0.31\linewidth}|Y|}
\hline
\textbf{Objeto} & \textbf{Autoria e campos mínimos} & \textbf{Invariantes} \\
\hline
\texttt{NodeContext} &
runtime: pedido, objetivo, \texttt{doneWhen}, estado projetado, skills e tools &
ancestrais concluídos; critérios, contagens e somente observações citadas;
sem \texttt{callId}; bodies delimitados; menu igual a \(T(g)\) \\
\hline
\texttt{NodeDecision} &
modelo: status, critérios com \texttt{criterionIndex}, satisfação, evidência em prosa e
\texttt{observationIndices[]}, resultado, pedido de revisão e razão &
índices zero-based, únicos, crescentes e mínimos; não contém observações, \texttt{callId}
nem referências opacas do provedor \\
\hline
\texttt{Observation} &
runtime: \texttt{goalId}, \texttt{toolName}, \texttt{callId}, entrada, saída e ordem &
chamada e retorno correlacionados exatamente uma vez; \texttt{callId} permanece interno;
a tool terminal de saída estruturada não é observação \\
\hline
\texttt{NodeOutcome} &
runtime: todos os campos de \texttt{NodeDecision} mais \texttt{observations[]} &
contém cada retorno bem-sucedido de tool executável, na ordem dos resultados;
referências fora do ledger invalidam o outcome antes de qualquer promoção \\
\hline
\texttt{RevisionRequest} &
modelo: \texttt{goalId}, \texttt{invalidatedAssumption}, \texttt{requestedEffect} &
\texttt{needs\_revision} exige pelo menos uma observação; o runtime associa
automaticamente o conjunto completo produzido pelo nó \\
\hline
\texttt{Artifact} &
runtime/modelo: \texttt{kind}, \texttt{mime} e \texttt{data} ou \texttt{reference} &
união estrita entre conteúdo inline e referência opaca; ordem preservada;
resolução e persistência fora do núcleo \\
\hline
\texttt{FinalDelivery} &
runtime: objetivos terminais, partes e artefatos discriminados &
construída sem nova chamada ao modelo; ordem estável e rastreável \\
\hline
\end{tabularx}
\par\smallskip\raggedright\small Fonte: elaboração própria (2026).
\end{table}
